Im ersten Schritt (isotherm) gilt $p_1V_1 = p_2V_2$ mit $V_2 = \left(1+\dVrO\right)V_1$:
\begin{align}
 \frac{p_2}{p_1} &= \pqaF = \frac{V_1}{\left(1+\dVrO\right)V_1} = \pqa
\end{align}
Im zweiten Schritt (isochor) gilt $p_2/T_2 = p_3/T_3$ mit $T_3 = \left(1-\dTrO\right)T_2$:
\begin{align}
 \frac{p_3}{p_2} &= \pqbF = \pqb
\end{align}
Insgesamt ist
\begin{align}
 \frac{p_3}{p_1} &= \pqF = \pqa\cdot\pqb = \pq
\end{align}
Der Druck ändert sich also um
\begin{align}
 \frac{\Delta p}{p_1} &= \dprF = \dpr \approx \result{\dprP{0}{2}}
\end{align}
er nimmt um rund einen Viertel ab.
